\documentclass[reqno]{amsart}
\usepackage[T1]{fontenc}
\usepackage{amsmath,amssymb,amsthm}
\usepackage{xr-hyper}
\usepackage[hidelinks]{hyperref}
\externaldocument[st-]{20261124T214500--supersingular-theta-divisibility__workbook}
\externaldocument[pf-]{20261125T013700--pairings-on-a-formal-group-of-height-two__workbook}

% Chapter number; \ref{chap} gives it.
\newcounter{chap}
% Chapter numbers, generated from 20261125T012000--grh-workbook__guide.org
% by scripts/make-chapters.py.  Do not edit: change the order of the
% index and rerun the script (./preview-site.sh and CI do this).
%
% Usage in a sheet, after \newcounter{chap}:
%   % Chapter numbers, generated from 20261125T012000--grh-workbook__guide.org
% by scripts/make-chapters.py.  Do not edit: change the order of the
% index and rerun the script (./preview-site.sh and CI do this).
%
% Usage in a sheet, after \newcounter{chap}:
%   % Chapter numbers, generated from 20261125T012000--grh-workbook__guide.org
% by scripts/make-chapters.py.  Do not edit: change the order of the
% index and rerun the script (./preview-site.sh and CI do this).
%
% Usage in a sheet, after \newcounter{chap}:
%   \input{chapters}\chapnumber{<sheet ID>}
\makeatletter
\providecommand\chapters@set[2]{\@namedef{chapters@#1}{#2}}
\chapters@set{20261125T001000}{1}% zeros-to-ramp-sums
\chapters@set{20261125T001100}{2}% inert-primes-and-estimate-b
\chapters@set{20261125T005000}{3}% galois-fourier-and-averaging
\chapters@set{20261125T005500}{4}% columns-from-a-sphere
\chapters@set{20261125T000000}{5}% slope-method
\chapters@set{20261126T030100}{6}% a-twist-for-each-prime
\chapters@set{20261125T003000}{7}% the-chain-and-the-square-root
\chapters@set{20261125T010000}{8}% rank-lemma-and-zero-estimate
\chapters@set{20261125T010500}{9}% poincare-series-and-near-orthogonality
\chapters@set{20261124T223000}{10}% elliptic-curves-for-the-proof
\chapters@set{20261125T009500}{11}% fock-model-and-theta-intertwiners
\chapters@set{20261124T223100}{12}% complex-multiplication-and-inert-reduction
\chapters@set{20261125T008000}{13}% one-cm-determinant
\chapters@set{20261124T214500}{14}% supersingular-theta-divisibility
\chapters@set{20261124T214600}{15}% the-isotropic-plane
\chapters@set{20261125T013700}{16}% pairings-on-a-formal-group-of-height-two
\chapters@set{20261125T013500}{17}% hecke-nondegeneracy-on-the-tree
\chapters@set{20261126T030200}{18}% the-bound-is-attained
\chapters@set{20261125T009000}{19}% kronecker-limit-and-chowla-selberg
\chapters@set{20261125T002000}{20}% definite-theta-lifts
\chapters@set{20261125T015500}{21}% siegel-weil-in-the-smallest-case
\chapters@set{20261125T002100}{22}% definite-theta-lifts-and-norms
\chapters@set{20261125T016500}{23}% hecke-operators-and-theta-lifts
\chapters@set{20261125T016700}{24}% divisibility-along-the-chain
\chapters@set{20261125T004000}{25}% denominators-and-sha
\chapters@set{20261125T004100}{26}% tate-twists-and-sha
\chapters@set{20261126T020430}{27}% the-minus-part-by-gauss-sums
\chapters@set{20261125T014100}{28}% congruence-modules
\chapters@set{20261125T014500}{29}% eisenstein-congruences-and-selmer-classes
\chapters@set{20261125T014700}{30}% pseudocharacters-and-the-length-bound
\chapters@set{20261125T011000}{31}% galois-representations-black-box
\chapters@set{20261126T030300}{32}% the-data
\chapters@set{20261126T020210}{33}% bessel-poincare-series-and-their-gram-operators
\chapters@set{20261126T030400}{34}% the-seesaw-identity
\chapters@set{20261126T030500}{35}% divisibility-of-the-composite-lift
\chapters@set{20261126T030600}{36}% theorem-a-short-forms
\chapters@set{20261126T030700}{37}% the-determinant-of-a-chain
\chapters@set{20261126T030800}{38}% the-cost-of-one-lift
\chapters@set{20261126T030900}{39}% denominators-of-projections-to-lifts
\chapters@set{20261126T031000}{40}% twists-with-few-congruences
\chapters@set{20261126T031100}{41}% theorem-c-and-the-main-inequality
\providecommand\chapnumber[1]{%
  \@ifundefined{chapters@#1}%
    {\PackageError{chapters}{Sheet #1 is not listed in the index}%
      {List it in 20261125T012000--grh-workbook__guide.org\MessageBreak
       and run scripts/make-chapters.py.}}%
    {\edef\thechap{\@nameuse{chapters@#1}}}}
\makeatother
\chapnumber{<sheet ID>}
\makeatletter
\providecommand\chapters@set[2]{\@namedef{chapters@#1}{#2}}
\chapters@set{20261125T001000}{1}% zeros-to-ramp-sums
\chapters@set{20261125T001100}{2}% inert-primes-and-estimate-b
\chapters@set{20261125T005000}{3}% galois-fourier-and-averaging
\chapters@set{20261125T005500}{4}% columns-from-a-sphere
\chapters@set{20261125T000000}{5}% slope-method
\chapters@set{20261126T030100}{6}% a-twist-for-each-prime
\chapters@set{20261125T003000}{7}% the-chain-and-the-square-root
\chapters@set{20261125T010000}{8}% rank-lemma-and-zero-estimate
\chapters@set{20261125T010500}{9}% poincare-series-and-near-orthogonality
\chapters@set{20261124T223000}{10}% elliptic-curves-for-the-proof
\chapters@set{20261125T009500}{11}% fock-model-and-theta-intertwiners
\chapters@set{20261124T223100}{12}% complex-multiplication-and-inert-reduction
\chapters@set{20261125T008000}{13}% one-cm-determinant
\chapters@set{20261124T214500}{14}% supersingular-theta-divisibility
\chapters@set{20261124T214600}{15}% the-isotropic-plane
\chapters@set{20261125T013700}{16}% pairings-on-a-formal-group-of-height-two
\chapters@set{20261125T013500}{17}% hecke-nondegeneracy-on-the-tree
\chapters@set{20261126T030200}{18}% the-bound-is-attained
\chapters@set{20261125T009000}{19}% kronecker-limit-and-chowla-selberg
\chapters@set{20261125T002000}{20}% definite-theta-lifts
\chapters@set{20261125T015500}{21}% siegel-weil-in-the-smallest-case
\chapters@set{20261125T002100}{22}% definite-theta-lifts-and-norms
\chapters@set{20261125T016500}{23}% hecke-operators-and-theta-lifts
\chapters@set{20261125T016700}{24}% divisibility-along-the-chain
\chapters@set{20261125T004000}{25}% denominators-and-sha
\chapters@set{20261125T004100}{26}% tate-twists-and-sha
\chapters@set{20261126T020430}{27}% the-minus-part-by-gauss-sums
\chapters@set{20261125T014100}{28}% congruence-modules
\chapters@set{20261125T014500}{29}% eisenstein-congruences-and-selmer-classes
\chapters@set{20261125T014700}{30}% pseudocharacters-and-the-length-bound
\chapters@set{20261125T011000}{31}% galois-representations-black-box
\chapters@set{20261126T030300}{32}% the-data
\chapters@set{20261126T020210}{33}% bessel-poincare-series-and-their-gram-operators
\chapters@set{20261126T030400}{34}% the-seesaw-identity
\chapters@set{20261126T030500}{35}% divisibility-of-the-composite-lift
\chapters@set{20261126T030600}{36}% theorem-a-short-forms
\chapters@set{20261126T030700}{37}% the-determinant-of-a-chain
\chapters@set{20261126T030800}{38}% the-cost-of-one-lift
\chapters@set{20261126T030900}{39}% denominators-of-projections-to-lifts
\chapters@set{20261126T031000}{40}% twists-with-few-congruences
\chapters@set{20261126T031100}{41}% theorem-c-and-the-main-inequality
\providecommand\chapnumber[1]{%
  \@ifundefined{chapters@#1}%
    {\PackageError{chapters}{Sheet #1 is not listed in the index}%
      {List it in 20261125T012000--grh-workbook__guide.org\MessageBreak
       and run scripts/make-chapters.py.}}%
    {\edef\thechap{\@nameuse{chapters@#1}}}}
\makeatother
\chapnumber{<sheet ID>}
\makeatletter
\providecommand\chapters@set[2]{\@namedef{chapters@#1}{#2}}
\chapters@set{20261125T001000}{1}% zeros-to-ramp-sums
\chapters@set{20261125T001100}{2}% inert-primes-and-estimate-b
\chapters@set{20261125T005000}{3}% galois-fourier-and-averaging
\chapters@set{20261125T005500}{4}% columns-from-a-sphere
\chapters@set{20261125T000000}{5}% slope-method
\chapters@set{20261126T030100}{6}% a-twist-for-each-prime
\chapters@set{20261125T003000}{7}% the-chain-and-the-square-root
\chapters@set{20261125T010000}{8}% rank-lemma-and-zero-estimate
\chapters@set{20261125T010500}{9}% poincare-series-and-near-orthogonality
\chapters@set{20261124T223000}{10}% elliptic-curves-for-the-proof
\chapters@set{20261125T009500}{11}% fock-model-and-theta-intertwiners
\chapters@set{20261124T223100}{12}% complex-multiplication-and-inert-reduction
\chapters@set{20261125T008000}{13}% one-cm-determinant
\chapters@set{20261124T214500}{14}% supersingular-theta-divisibility
\chapters@set{20261124T214600}{15}% the-isotropic-plane
\chapters@set{20261125T013700}{16}% pairings-on-a-formal-group-of-height-two
\chapters@set{20261125T013500}{17}% hecke-nondegeneracy-on-the-tree
\chapters@set{20261126T030200}{18}% the-bound-is-attained
\chapters@set{20261125T009000}{19}% kronecker-limit-and-chowla-selberg
\chapters@set{20261125T002000}{20}% definite-theta-lifts
\chapters@set{20261125T015500}{21}% siegel-weil-in-the-smallest-case
\chapters@set{20261125T002100}{22}% definite-theta-lifts-and-norms
\chapters@set{20261125T016500}{23}% hecke-operators-and-theta-lifts
\chapters@set{20261125T016700}{24}% divisibility-along-the-chain
\chapters@set{20261125T004000}{25}% denominators-and-sha
\chapters@set{20261125T004100}{26}% tate-twists-and-sha
\chapters@set{20261126T020430}{27}% the-minus-part-by-gauss-sums
\chapters@set{20261125T014100}{28}% congruence-modules
\chapters@set{20261125T014500}{29}% eisenstein-congruences-and-selmer-classes
\chapters@set{20261125T014700}{30}% pseudocharacters-and-the-length-bound
\chapters@set{20261125T011000}{31}% galois-representations-black-box
\chapters@set{20261126T030300}{32}% the-data
\chapters@set{20261126T020210}{33}% bessel-poincare-series-and-their-gram-operators
\chapters@set{20261126T030400}{34}% the-seesaw-identity
\chapters@set{20261126T030500}{35}% divisibility-of-the-composite-lift
\chapters@set{20261126T030600}{36}% theorem-a-short-forms
\chapters@set{20261126T030700}{37}% the-determinant-of-a-chain
\chapters@set{20261126T030800}{38}% the-cost-of-one-lift
\chapters@set{20261126T030900}{39}% denominators-of-projections-to-lifts
\chapters@set{20261126T031000}{40}% twists-with-few-congruences
\chapters@set{20261126T031100}{41}% theorem-c-and-the-main-inequality
\providecommand\chapnumber[1]{%
  \@ifundefined{chapters@#1}%
    {\PackageError{chapters}{Sheet #1 is not listed in the index}%
      {List it in 20261125T012000--grh-workbook__guide.org\MessageBreak
       and run scripts/make-chapters.py.}}%
    {\edef\thechap{\@nameuse{chapters@#1}}}}
\makeatother
\chapnumber{20261124T214600}
\renewcommand{\thesection}{\thechap.\arabic{section}}

\theoremstyle{definition}
\newtheorem{problem}{Problem}[section]
\theoremstyle{plain}
\newtheorem{theorem}{Theorem}[section]

\newcommand{\Z}{\mathbb{Z}}
\newcommand{\Q}{\mathbb{Q}}
\newcommand{\C}{\mathbb{C}}
\newcommand{\F}{\mathbb{F}}
\newcommand{\ii}{\mathrm{i}}
\newcommand{\Oc}{\mathcal{O}}
\newcommand{\sll}{\operatorname{sl}}
\newcommand{\Frob}{\operatorname{Frob}}

\setlength{\emergencystretch}{2em}

\title{The isotropic plane}
\author{}
\date{November 24, 2026}

\begin{document}
\maketitle
\refstepcounter{chap}\label{chap}

\paragraph{Prerequisites.}
Chapter~\ref{st-chap}, \emph{The invariant derivation at a supersingular prime}: the moments \(t(a,w)\), the functions \(F\) and \(\Theta\) and the directions \(e_\pm\) of Section~\ref{st-sec:moments}, the formal group \(\hat E\) and the ring \(\Oc = \Z_p[\ii]\) of Section~\ref{st-sec:lemniscate}, the function \(E_p\), Problem~\ref{st-pr:first-gain} and Theorem~\ref{st-thm:single}, and Section~\ref{st-sec:one-theta}.  Quasi-periods of \(\sigma\) and Legendre's relation \cite[\S I.5]{Si2}.  Decomposition groups in cyclotomic fields.  Rational functions and divisors on an abelian scheme over a discrete valuation ring, at the level of the statement that a rational function whose divisor avoids a point is regular there \cite[Lemma~2.13]{BK}.

\section{The isotropic plane}\label{sec:plane}

The function \(F\) does not gain by itself (Problem~\ref{st-pr:failure}).  Why should \(\Theta\) along the directions \(e_\pm = (1, \pm\zeta)\) do better?  They are the eigenlines of an action of \(\Z[\zeta]\) on \(\Z[\ii]^2\).  Let
\[
  \rho(\zeta) := \begin{pmatrix} 0 & 1 \\ \ii & 0 \end{pmatrix},
\]
let \(\sigma_+\) be the identity of \(\Q(\zeta)\) and \(\sigma_-\) the automorphism \(\zeta \mapsto -\zeta\), and let \(H(Z, W) := Z_1 \bar W_1 + Z_2 \bar W_2\).

\begin{problem}\label{pr:rho}
\(\rho\) extends to a ring homomorphism \(\Z[\zeta] \to M_2(\Z[\ii])\), and \(\rho(\bar a) = \rho(a)^*\), the conjugate transpose.
\end{problem}

\begin{problem}\label{pr:eigen}
\(\rho(a) e_\pm = \sigma_\pm(a) e_\pm\), and \(H(e_+, e_-) = 0\).
\end{problem}

For each prime \(p \equiv 3 \pmod 4\), fix an embedding \(\Q(\zeta) \hookrightarrow \Q_p(\ii)\) extending the embedding of \(\Q(\ii)\) given by \(\Oc = \Z_p[\ii]\), and use it for \(v_p\) and for reduction modulo \(p\).  Let \(\sigma_p\) be the nonidentity automorphism of \(\Q_p(\ii)/\Q_p\).

Now let \(p \equiv 3 \pmod 8\).  The Frobenius \(\zeta \mapsto \zeta^3\) generates the decomposition group of \(p\) in \(\Q(\zeta)\), with fixed field \(\Q(\sqrt{-2})\), \(\sqrt{-2} := \zeta + \zeta^3\).  Since \(\Z[\sqrt{-2}]\) is a principal ideal domain, \(p = x^2 + 2y^2\) for integers \(x\), \(y\).

\begin{problem}\label{pr:a}
After replacing \(y\) by \(-y\) if necessary, \(a := x + y\sqrt{-2}\) lies in \(p\Oc\), while \(\bar a = x - y\sqrt{-2}\) is a unit of \(\Oc\).  Moreover \(\sigma_-(a) = \bar a\).
\end{problem}

\begin{problem}\label{pr:a-fails}
If \(p \equiv 7 \pmod 8\), then \(v_p(a) = v_p(\bar a)\) for every \(a \in \Z[\zeta]\).
\end{problem}

On Gaussian coefficients \(\sigma_p\) agrees with complex conjugation.  The local form \(H_p(z, w) := z_1 \sigma_p(w_1) + z_2 \sigma_p(w_2)\) on \(\Oc^2\) reduces to the hermitian form \(h(z, w) := z_1 w_1^p + z_2 w_2^p\) on \(\F_{p^2}^2\).  Let \(\bar e_\pm\) be the reductions of \(e_\pm \in \Oc^2\).  Write \(\widetilde{\zeta}\) for the reduction of \(\zeta\) modulo \(p\).

\begin{problem}\label{pr:isotropic}
\(h(\bar e_+, \bar e_+) = h(\bar e_-, \bar e_-) = 1 + \widetilde{\zeta}^{p+1}\) and \(h(\bar e_+, \bar e_-) = 1 - \widetilde{\zeta}^{p+1}\).  So for \(p \equiv 3 \pmod 8\) both lines are isotropic and \(h\) pairs them, while for \(p \equiv 7 \pmod 8\) they are anisotropic and orthogonal.
\end{problem}

We call \(e_-\) the \emph{partner} of \(e_+\).  In terms of embeddings, \(\sigma_- = c \circ \Frob \circ \sigma_+\), where \(c\) is complex conjugation and \(\Frob(\zeta) = \zeta^p\); for \(p \equiv 7 \pmod 8\), by contrast, \(c \circ \Frob \circ \sigma_+ = \sigma_+\).  The derivatives along \(e_+\) will gain, and those along the partner \(e_-\) will pay for it.

\section{Integrality on an eigenline}\label{sec:line}

Let \(P(s) := \Theta(s e_+) = F(s)\, F(\zeta s)\).  Since \(F\) is even, \(\Theta(t e_-) = P(t)\), and \(t(n, 0) = P^{(n)}(0)\).  Write \(P(s) = \hat P(\sll s)\) with \(\hat P \in \Q[[X]]\).  We show that \(\hat P\) is \(p\)-integral for every prime \(p \equiv 3 \pmod 8\), unlike \(F\).

Let \(L\) be the period lattice of \(\sll\) and \(\eta\) the quasi-period map of \(\Lambda\).  Since \(\eta(\ii\omega) = -\ii\eta(\omega)\), Legendre's relation gives \(\eta(\omega) = (\pi/\operatorname{vol}\Lambda)\, \bar\omega\), and with the law \(\sigma(u + \omega) = \pm e^{\eta(\omega)(u + \omega/2)}\sigma(u)\) \cite[\S I.5]{Si2},
\[
  F(u + \omega) = \exp\bigl(c\, \bar\omega u + \tfrac{c}{2} \lvert\omega\rvert^2\bigr) F(u),
  \qquad \omega \in \Lambda,\quad c := \frac{2\pi}{\operatorname{vol}\Lambda} .
\]

\begin{problem}\label{pr:Fa}
For \(\gamma \in L^2\), \(\Theta(Z + \gamma) = \exp\bigl(c H(Z, \gamma) + \tfrac{c}{2}H(\gamma, \gamma)\bigr)\Theta(Z)\).  For \(a \in \Z[\zeta]\), \(a \ne 0\), the function
\[
  F_a(Z) := \frac{\Theta(\rho(a) Z)}{\Theta(\rho(\bar a) Z)}
\]
is \(L^2\)-periodic with \(F_a(0) = 1\).  It descends to a rational function on \(E \times E\) defined over \(\Q(\ii)\).
\end{problem}

\begin{problem}\label{pr:zeros}
On \(\C/L = E(\C)\), the zeros of \(F\) are the two points at infinity of \(y^2 = 1 - X^4\), each with multiplicity two.
\end{problem}

Let \(\mathcal{A}\) be \(E \times E\) over \(\Oc\), an abelian scheme on which \(\Z[\zeta]\) acts through \(\rho\), and write \(\hat F_a(X_1, X_2)\) for the expansion of \(F_a\) in \(X_1 = \sll Z_1\), \(X_2 = \sll Z_2\).

\begin{problem}\label{pr:Fa-integral}
\(\hat F_a \in \Oc[[X_1, X_2]]\).
\end{problem}

Now let \(p \equiv 3 \pmod 8\), and let \(a\) be as in Problem~\ref{pr:a}.

\begin{problem}\label{pr:functional}
\(\hat P([a]X) = \hat F_a(X, [\zeta]X)\, \hat P([\bar a]X)\).
\end{problem}

\begin{problem}\label{pr:BK-induction}
The coefficients of \(\hat P\) lie in \(\Z_{(p)}\).
\end{problem}

Problems~\ref{pr:Fa}--\ref{pr:BK-induction} and Theorem~\ref{st-thm:single} prove the following.

\begin{theorem}\label{thm:line}
Let \(p \equiv 3 \pmod 8\).  Then \(\hat P \in \Z_{(p)}[[X]]\), and \(v_p(t(n, 0)) = v_p(P^{(n)}(0)) \ge E_p(n)\) for every \(n \ge 0\).
\end{theorem}

The induction in Problem~\ref{pr:BK-induction} is the argument Bannai and Kobayashi use for the Kronecker theta function at ordinary primes \cite[Prop.~2.14]{BK}, where one element of the CM order serves for all coordinates at once.  Here \(\mathcal{A}\) is supersingular, but on the line \(\C e_+\) the element \(a\) acts by a non-unit and \(\bar a\) by a unit.  At \(p \equiv 7 \pmod 8\) there is no such element (Problem~\ref{pr:a-fails}), and at \(p = 7\) the conclusion fails.

\begin{problem}\label{pr:p7}
\(\hat P(X) = 1 - \tfrac{1}{70}X^8 + O(X^{12})\).  So \(\hat P\) is not \(7\)-integral, and \(t(8,0) = -576\).
\end{problem}

\section{The cross term}\label{sec:cross}

Theorem~\ref{thm:line} accounts for the derivatives along each axis.  What remains is a factor equal to \(1\) on both axes.  Throughout, \(p \equiv 3 \pmod 8\), \(X_1 = \sll s\), \(X_2 = \sll t\), and \(P_n := P^{(n)}(0)\).

\begin{problem}\label{pr:J}
\(J(s,t) := \partial_s\partial_t \log\Theta(s e_+ + t e_-) = 2\sll^2(s + t) - 2\ii\, \sll^2\bigl(\zeta(s - t)\bigr)\), and \(J \in \Oc[[X_1, X_2]]\).
\end{problem}

Let \(M(s,t) := \int_0^s \int_0^t J\), and \(M_{\alpha\beta} := \partial_s^\alpha \partial_t^\beta M(0,0)\).

\begin{problem}\label{pr:factor}
\(\Theta(s e_+ + t e_-) = P(s)\, P(t)\, e^{M(s,t)}\).  The numbers \(M_{\alpha\beta}\) vanish unless \(\alpha, \beta \ge 1\), and then \(M_{\alpha\beta} \in p^{E_p(\alpha - 1)}\Oc\).
\end{problem}

\begin{problem}\label{pr:expformula}
Let \(S\) and \(T\) be disjoint sets with \(a\) and \(w\) elements, respectively.  Then
\[
  t(a,w) = \sum P_{\lvert S_0 \rvert}\, P_{\lvert T_0 \rvert} \prod_{B \in \pi} M_{\lvert B \cap S \rvert, \lvert B \cap T \rvert},
\]
the sum over subsets \(S_0 \subseteq S\), \(T_0 \subseteq T\), and partitions \(\pi\) of \((S \smallsetminus S_0) \sqcup (T \smallsetminus T_0)\) into blocks each meeting both \(S\) and \(T\).
\end{problem}

\begin{problem}\label{pr:first-level}
Each term in Problem~\ref{pr:expformula} has at most \(w\) blocks and valuation at least \(\lfloor a/p \rfloor - w\).
\end{problem}

Thus \(v_p(t(a,w)) \ge \lfloor a/p \rfloor - w\).  The \emph{type} of a term in Problem~\ref{pr:expformula} consists of \(a_0 = \lvert S_0 \rvert\), \(w_0 = \lvert T_0 \rvert\), and the numbers \(m_{\alpha\beta}\) of blocks with \(\lvert B \cap S \rvert = \alpha\) and \(\lvert B \cap T \rvert = \beta\).

\begin{problem}\label{pr:multiplicity}
For every type \((a_0, w_0, (m_{\alpha\beta}))\) occurring in Problem~\ref{pr:expformula}, the number of terms of that type is \(N \cdot a!/\bigl(a_0! \prod_{\alpha,\beta} (\alpha!)^{m_{\alpha\beta}}\bigr)\) for some positive integer \(N\).
\end{problem}

\begin{problem}\label{pr:levels}
Let \(\alpha_1, \dotsc, \alpha_k\) be positive integers and \(a_0\) a nonnegative integer such that \(a_0 + \sum_{b=1}^k \alpha_b = a\).  Let \(n_i := \#\{b : p^i \mid \alpha_b\}\).  Then
\[
  v_p\Bigl(\frac{a!}{a_0!\, \alpha_1! \dotsm \alpha_k!}\Bigr) + E_p(a_0) + \sum_{b=1}^k E_p(\alpha_b - 1)
  \ \ge\ \sum_{\substack{i \ge 1 \\ i \text{ odd}}} \Bigl(\Bigl\lfloor \frac{a}{p^i} \Bigr\rfloor - n_i\Bigr)_+ ,
\]
where \(x_+ := \max(x, 0)\).
\end{problem}

Problems~\ref{pr:factor}--\ref{pr:levels}, Theorem~\ref{thm:line} and the symmetry \(t(a,w) = t(w,a)\) of Problem~\ref{st-pr:t-sym} give the following theorem.

\begin{theorem}\label{thm:main}
Let \(p \equiv 3 \pmod 8\).  For all integers \(a, w \ge 0\),
\[
  v_p\bigl(t(a,w)\bigr) \ \ge\ \max\Biggl( \sum_{\substack{i \ge 1 \\ i \text{ odd}}} \Bigl(\Bigl\lfloor \frac{a}{p^i} \Bigr\rfloor - w\Bigr)_+ ,\ \sum_{\substack{i \ge 1 \\ i \text{ odd}}} \Bigl(\Bigl\lfloor \frac{w}{p^i} \Bigr\rfloor - a\Bigr)_+ \Biggr) .
\]
\end{theorem}

In the table of Section~\ref{st-sec:moments} the bound is an equality whenever it is positive.  It is not always sharp: \(t(11,1) = -2^9 \cdot 3^6\), while the bound at \(p = 3\) is \(2\).

\subsection*{Exercise}
Let \(M = s^p t/p\) and let \(w \ge 0\) be an integer.  Show that \(\partial_s^{pw}\partial_t^w e^{M}(0,0) = (pw)!/p^w\), with valuation \(v_p(w!)\).  Deduce that this valuation is \(0\) for \(0 \le w < p\).  Show also that for every \(G \in \Oc[[X_1, X_2]]\), with \(X_1 = \sll s\) and \(X_2 = \sll t\), the valuation of \(\partial_s^{pw}\partial_t^w G(0,0)\) is at least \(w\).

\section{Questions}\label{sec:questions}

Theorem~\ref{thm:main} does not extend to \(p \equiv 7 \pmod 8\): it fails at \(p = 7\) (Problem~\ref{pr:p7}).  What fails is the choice of directions: \(e_\pm\) are global eigenlines of \(\Z[\zeta]\), and they are isotropic modulo \(p\) only for \(p \equiv 3 \pmod 8\) (Problem~\ref{pr:isotropic}).  For every prime \(p \equiv 3 \pmod 4\), the local form \(H_p\) on \(\Oc^2\) admits a basis \(f, f'\) with \(H_p(f, f) = H_p(f', f') = 0\) and \(H_p(f, f') = 1\).  If \(p \equiv 7 \pmod 8\), neither line \(\Q_p(\ii)e_\pm\) contains a nonzero isotropic vector.  Theorem~\ref{thm:main} concerns the two global eigenlines.

Let \(A\) be an abelian variety with complex multiplication by the ring of integers of a CM field \(K\), with CM type \(\Phi\), and let \(\Theta\) be a normalized theta function on its universal cover.  Fix a prime \(p\) and an embedding \(\jmath : \overline{\Q} \hookrightarrow \overline{\Q}_p\), so that \(\jmath \circ \phi\) determines a prime of \(K\) above \(p\).  The eigenlines \(\C e_\phi\) (\(\phi \in \Phi\)) replace \(e_\pm\); the case above is \(K = \Q(\zeta)\), \(\Phi = \{\sigma_+, \sigma_-\}\).  The argument of Section~\ref{sec:line} on \(\C e_\phi\) needs an element \(a\) of the CM order with \(\phi(a)\) a non-unit and \(\phi(\bar a)\) a unit, and such an element exists exactly when complex conjugation does not fix the prime of \(K\) determined by \(\phi\).  In the first sum of Theorem~\ref{thm:main}, the subtracted order is \(w\), the order along \(e_-\).  For \(p \equiv 3 \pmod 8\), this is the partner direction at every odd level \(i\), since \(c \circ \Frob^i \circ \sigma_+ = \sigma_-\).  For a general CM type, suppose \(p\) is unramified in the normal closure of \(K\), and let \(\Frob\) be the arithmetic Frobenius at the chosen prime of that closure.  At a level \(i\) with \(\Frob^i \circ \phi \notin \Phi\), is the order subtracted from the order along \(e_\phi\) the order along \(e_{c \circ \Frob^i \circ \phi}\)?

Among the odd supersingular primes for \(E\), the choice of a global CM element in Problem~\ref{pr:a} fails exactly for \(p \equiv 7 \pmod 8\).  What derivative estimates hold in local hyperbolic coordinates at these primes?  Let \(A_5\) be the Jacobian of \(y^2 = x^5 - 1\), with CM by \(\Q(\zeta_5)\).  Among odd primes \(p \ne 5\), the supersingular primes for \(A_5\) are exactly those with \(p \not\equiv 1 \pmod 5\), by the Shimura--Taniyama formula for the slopes of a Jacobian with complex multiplication \cite{LMPT}.  At these primes complex conjugation is in the decomposition group of every prime above \(p\), so there is no element of the CM order separating a prime from its conjugate.  For this Jacobian at such a prime, what choices of normalized theta function, eigenline direction and formal parameter from an integral model give a restriction with \(p\)-integral coefficients?  What valuation bound then holds for the derivatives in that direction at the origin?

\paragraph{Which step this chapter serves.}
This chapter gives a two-variable model for the partner terms in the divisibility estimate of the proof.  Theorem~\ref{pf-thm:count} of Chapter~\ref{pf-chap}, \emph{Pairings on a formal group of height two}, bounds the valuation of jets along products of minors for series in the prescribed normal form and at the prescribed scale.  Write \(r_e(\tau)\) for the partner of the infinite place \(\tau\) at level \(e\), as defined there.  Among the columns of height less than the number of rows, let \(a_\tau\) count the columns of height one at \(\tau\), and let \(a'_\tau\) count the other columns there.  In the positive-level sum for the predicted divisibility, the summand at \(\tau\) and \(e\) contains the positive part \(\bigl(a_\tau/p^e - a_{r_e(\tau)} - a'_{r_e(\tau)}\bigr)_+\).  The sum runs over all infinite places and over odd positive \(e\) if \(p\) is inert in \(\Q(\sqrt{-d})\), or all positive \(e\) if \(p = d\).  Its weight is \(1\) if \(p\) is inert and \(1/2\) if \(p = d\).  The order \(w\) along \(e_-\) in the first sum of Theorem~\ref{thm:main} plays the analogous role.  Theorem~\ref{pf-thm:count} assumes that the prime \(p\) does not split in \(\Q(\sqrt{-d})\), with no separate congruence condition on \(p\) modulo \(8\).  Theorems~\ref{thm:line} and~\ref{thm:main} give the estimates for the workbook moments.

\paragraph{In the Lean formalization.}
Theorem~\ref{pf-thm:count} is the valuation form of \texttt{Parameters.norm\_jet\_minorProd\_le\_of\_hasNormalForm} in \texttt{grh/src/New/Derivatives/Divisibility/Count.lean} of the Lean proof \cite{GL}.  Its predicted divisibility is \texttt{predictedDivisibility} in the same file.  The positive partner levels are \texttt{positivePartnerExponents} in \texttt{grh/src/Lib/CM/FormalGroup/PrimeNotation.lean}, and the partner map is \texttt{Parameters.partner} in \texttt{grh/src/New/Derivatives/FormalGroup/ThetaProducts.lean}.

\section*{Hints}

\subsection*{\S\ref{sec:plane}}
\begin{description}
\item[\ref{pr:rho}] \(\rho(\zeta)^2 = \ii\), and \(\rho(\zeta)\) is unitary.
\item[\ref{pr:a}] \(a \bar a = p\), and \(\Oc\) is a discrete valuation ring with uniformizer \(p\).  For \(p = 3\), take \(x = 1\) and test the two signs of \(y = \pm 1\) in the chosen embedding.
\item[\ref{pr:a-fails}] The Frobenius automorphism of \(\Oc\) preserves \(v_p\) and maps \(\zeta\) to \(\zeta^p\).
\item[\ref{pr:isotropic}] Compare the actions of \(\sigma_p\) and of complex conjugation on \(\zeta\).
\end{description}

\subsection*{\S\ref{sec:line}}
\begin{description}
\item[\ref{pr:Fa}] \(\rho(a)\) preserves \(L^2\) because \(\ii L = L\), and \(\rho(a)^*\rho(a) = \rho(a\bar a) = \rho(\bar a)^*\rho(\bar a)\).  For the field of definition, compare the Taylor expansion at the origin with its transforms under automorphisms of \(\C\) over \(\Q(\ii)\).
\item[\ref{pr:zeros}] \(F = \sigma^2/\sll^2\) by Problem~\ref{st-pr:wp-sl}.
\item[\ref{pr:Fa-integral}] Use Problem~\ref{pr:zeros} to locate the horizontal components of the divisor near the origin of the special fiber.  How can the value \(F_a(0)\) determine the remaining vertical multiplicity?
\item[\ref{pr:functional}] Evaluate \(\Theta(\rho(a)Z) = F_a(Z)\Theta(\rho(\bar a)Z)\) at \(Z = s e_+\), using Problems~\ref{pr:eigen} and~\ref{st-pr:O-action}.
\item[\ref{pr:BK-induction}] Induction on the degree, comparing coefficients of \(X^n\) in Problem~\ref{pr:functional}.
\end{description}

\subsection*{\S\ref{sec:cross}}
\begin{description}
\item[\ref{pr:J}] Problem~\ref{st-pr:logF}.  Then \(\sll(s+t) = X_1 \oplus X_2\) and \(\sll(\zeta(s-t)) = [\zeta](X_1 \ominus X_2)\).
\item[\ref{pr:factor}] Restrict \(\log\Theta\) to the two axes.  For the bound, express \(M_{\alpha\beta}\) through derivatives of \(J\), and use Theorem~\ref{st-thm:single}.
\item[\ref{pr:expformula}] Expand \(e^M = \sum_k M^k/k!\) and distribute the \(a + w\) labelled derivatives.
\item[\ref{pr:first-level}] Each block meets \(T\).  Keep only the first level of each \(E_p\).
\item[\ref{pr:multiplicity}] Choose the \(T\)-parts first; after that the blocks are distinguishable.
\item[\ref{pr:levels}] Apply Legendre's formula level by level, and use \(\lfloor \alpha/p^i \rfloor - \lfloor (\alpha-1)/p^i \rfloor \in \{0, 1\}\).  Each level contributes a nonnegative amount.
\end{description}

\begin{thebibliography}{Si2}
\bibitem[BK]{BK} K.~Bannai and S.~Kobayashi, \emph{Algebraic theta functions and the \(p\)-adic interpolation of Eisenstein--Kronecker numbers}, Duke Math.\ J.\ \textbf{153} (2010), 229--295; arXiv:math/0610163.  Numbering as in the arXiv version.
\bibitem[GL]{GL} L.~Sernau, \emph{GRH for Dirichlet \(L\)-functions, in Lean~4}, a shorter rewriting of the Lean proof, text by Claude, \url{https://github.com/lsernau/grh-lean}, commit \texttt{72fead1}, readable edition \texttt{grh/}.
\bibitem[LMPT]{LMPT} W.~Li, E.~Mantovan, R.~Pries and Y.~Tang, \emph{Newton polygons of cyclic covers of the projective line branched at three points}, arXiv:1805.04598.
\bibitem[Si2]{Si2} J.~H.~Silverman, \emph{Advanced Topics in the Arithmetic of Elliptic Curves}, Grad.\ Texts in Math.\ 151, Springer, 1994.
\end{thebibliography}

\end{document}
